\chapter{Divergence of detachment forces in the finite Voronoi model}
\label{chap:voronoi}

\begin{singlespace}
    This chapter is adapted from the following preprint:
    \begin{itemize}
        \item \fullcite{wang2026divergence}.
    \end{itemize}
\end{singlespace}


\section{Introduction}
\label{sec:ch4:intro}

Both confluent tissues, in which cells densely pack and tile space, and nonconfluent cell populations, in which cells or clusters of cells are separated by gaps, are common in biology. Many systems transition between these two states.
A prominent example is epithelial-mesenchymal transition (EMT), in which epithelial tissues lose cohesion and give rise to dispersed, migratory cells or clusters~\cite{youssef2024epithelial, tripathi2020physics}.
Similarly, during cancer dissemination, tumor cell collectives can transition from a compact, confluent mass to a nonconfluent, invasive state~\cite{harper2016mechanism, cheung2016collective}.
Understanding such transitions is central to describing collective cell behavior, and these approaches require simulation methods that can correctly resolve the mechanics of both confluent and nonconfluent tissues. One popular and powerful approach for simulating confluent tissues is the active Voronoi model~\cite{bi2016motility}, in which cells have a polygonal shape given by the Voronoi tessellation of cell centers, and cells evolve under a mechanical energy that depends on this geometry. These models have been particularly successful in explaining the jamming transition of confluent monolayers, where cells tile space without gaps and neighbor exchanges occur through topological rearrangements~\cite{bi2015density, bi2016motility, chen2022activation}---but the commonly simulated active Voronoi model~\cite{bi2016motility} does not allow for simulating nonconfluent tissues. However, recent work has developed so-called ``finite Voronoi'' models, originating from the concept of free Dirichlet domains~\cite{graner1993can}, to extend the conventional Voronoi model to describe nonconfluent monolayers~\cite{bock2010generalized, schaller2005multicellular, teomy2018confluent, huang2023bridging}. Finite Voronoi models retain the advantages of Voronoi-based descriptions in characterizing cell shape, while remaining significantly less computationally demanding than approaches like phase-field~\cite{nonomura2012study, palmieri2015multiple, chiang2024multiphase, wang2025confinement}, subcellular element~\cite{sandersius2008modeling, sweet2011modelling, sandersius2011emergent}, or deformable polygon~\cite{boromand2018jamming, lv2024active, weng2022convergent} models, making them promising candidates for simulations of cohesive clusters, active fragmentation, and tissue rupture~\cite{huang2023bridging, teomy2018confluent}.

Here, we study the finite Voronoi model of Refs.~\cite{teomy2018confluent,huang2023bridging}, and demonstrate that in some circumstances the force required to detach two cells can diverge unphysically due to the imposed geometry of cell-cell contacts, preventing cells from rearranging. This is akin to---though more dramatic than---earlier results showing that the assumption of a strict Voronoi shape can eliminate the unjamming transition observed in vertex models~\cite{sussman2018no} and change the dynamics of heterotypic interfaces~\cite{lawson2024differences,yue2024scale}.

We introduce the active finite Voronoi model in Sec.~\ref{sec:ch4:model}. We then demonstrate that when simulating tissue fracture, the results systematically and unavoidably depend on the time step used in numerical simulation, with fracture vanishing at small time steps (Sec.~\ref{sec:ch4:timestep}).
We then show that this pathology is not a numerical accident but a property of the finite Voronoi geometry. In Sec.~\ref{sec:ch4:doublet}, we analyze the simplest case of cell doublet detachment and show that the detachment force will diverge if there is an interfacial tension that differs between cell-cell and cell-medium interfaces.
Based on this analysis, we propose a regularization to address this divergence in Sec.~\ref{sec:ch4:regularization}.
In Sec.~\ref{sec:ch4:calibration}, we compare the finite Voronoi model with a deformable polygon model that has the same energy function, but does not assume the strict Voronoi shape, and analyze two different calibration strategies in Secs.~\ref{sec:ch4:cal1} and \ref{sec:ch4:cal2}.


\section{Model}
\label{sec:ch4:model}

We consider $N$ cells in two dimensions, labeled by $i=1,\dots,N$, with cell centers at positions $\{\mathbf r_i\}$.
In a conventional Voronoi model for confluent monolayers~\cite{bi2016motility}, Voronoi tessellations tile all of space, and the Voronoi region of cell $i$ is the set of points closer to $\mathbf{r}_i$ than to any other center.
In the finite Voronoi (FV) model, we assume that the cell cannot extend further than a distance $\ell$ away from its center. In this way, the shape of a cell is obtained by truncating the Voronoi region at a maximum radius $\ell$ from the cell center~\cite{teomy2018confluent,huang2023bridging}.
More formally, the domain of cell $i$ is the intersection of its Voronoi region with the disk $\{\mathbf x:|\mathbf x-\mathbf r_i|\leqslant \ell\}$.
As a result, FV boundaries are composed of (i) straight Voronoi segments where two cells are in contact and (ii) circular arcs of radius $\ell$ where the cell is exposed to the surrounding medium (Fig.~\ref{fig:ch4:FV_model}).
This construction allows two cells to separate once their centers move sufficiently far apart, creating gaps between cells without introducing additional degrees of freedom.

\begin{figure}[t]
\centering
\includegraphics[width=0.7\textwidth]{chap4/FV_model.pdf}
\caption{Illustration of the finite Voronoi (FV) model. Cell center positions $\{\mathbf{r}_i\}$ are the orange points, and gray lines are the corresponding Voronoi diagram generated by $\{\mathbf{r}_i\}$ while the dashed lines represent rays of Voronoi edges that extend to infinity.
Empty circles represent vertices $\{\mathbf{h}_m\}$ in the FV model: Blue circles are the triple junction vertices connecting three cells, while purple circles are the outer vertices $\mathbf{h}^{\textrm{out}}$ connecting two cells.
Scale bar: maximum diameter $2\ell$.}
\label{fig:ch4:FV_model}
\end{figure}

The energy of the cells is defined similarly to that in the conventional Voronoi/vertex-based model: a quadratic elastic energy penalizing area deviation away from a preferred area, a quadratic energy penalizing cell perimeter, and line tensions that are potentially different between cell-cell and cell-medium interfaces~\cite{huang2023bridging,teomy2018confluent,chen2022activation}:
\begin{equation}
\label{eq:ch4:energy-full}
    E=\sum_i K_A(A_i-A_0)^2+K_PP_i^2 + \lambda^{(c)} P_i^{(c)}+\lambda^{(n)}P_i^{(n)},
\end{equation}
where $K_A$ and $K_P$ are elastic moduli for cell area $A_i$ and perimeter $P_i$, respectively, and $A_0$ is the preferred cell area. $\lambda^{(c)}$ and $\lambda^{(n)}$ are the cortical tensions for contacting edges (cell-cell) and non-contacting edges (cell-medium interface), and the total circumference $P_i=P_i^{(c)}+P_i^{(n)}$ is the sum of contacting length and non-contacting length.
Up to an additive constant, this can be rewritten in a ``preferred perimeter'' form:
\begin{equation}
\label{eq:ch4:energy}
    E=\sum_i K_A(A_i-A_0)^2 + K_P(P_i-P_0)^2 + \Lambda P_i^{(n)},
\end{equation}
where $P_0=-\lambda^{(c)}/2 K_P$ is the preferred perimeter, and $\Lambda \equiv \lambda^{(n)}-\lambda^{(c)}$ measures the tension difference between contacting and non-contacting edges.

We assume that the cell centers are overdamped, so cell $i$ has one term in its velocity proportional to the force $-\nabla_i E$ on cell $i$ and one arising from self-propulsion~\cite{huang2023bridging},
\begin{equation}
\label{eq:ch4:EOM}
    \dot{\mathbf{r}}_i = -\mu\bm{\nabla}_i E + v_0\mathbf{n}_i,
\end{equation}
where $\mu$ is the cell mobility, $v_0$ sets the self-propulsion speed, and $\mathbf{n}_i=(\cos\theta_i,\sin\theta_i)$ is the polarity direction---the direction the cell would travel in the absence of cell-cell interactions. The polarity undergoes rotational diffusion,
\begin{equation}
\label{eq:ch4:diffusion}
    \dot{\theta}_i=\sqrt{2D_r}\eta_i(t),
\end{equation}
with zero-mean, unit-variance Gaussian white noise satisfying $\langle\eta_i(t)\rangle=0$ and $\langle\eta_i(t)\eta_j(t')\rangle=\delta_{ij}\delta(t-t')$.

All equations are nondimensionalized using $K_A$ and a characteristic length scale $L$, e.g., $K_P$ is scaled by $K_A L^2$ and $\Lambda$ is scaled by $K_A L^3$~\footnote{This corrects a typo in Eq.~(E4) of Ref.~\cite{teomy2018confluent}, where $\Lambda$ should be scaled by $K_AL^3$ rather than $K_AL^4$.}.
Following Ref.~\cite{huang2023bridging}, we choose the length scale $L=\sqrt{A_0/\pi}$ so that the nondimensional preferred area is $A_0=\pi$.
Time is scaled by the mechanical relaxation time $1/\mu K_A L^2$. Unless otherwise stated, we set the scaled maximum radius $\ell=1$ and the scaled $K_P=1$ (see Table~\ref{tab:ch4:params}), and we vary the control parameters $P_0$ and $\Lambda$ to explore the detachment forces. Though we have implemented the model independently, our model assumptions are identical to those of Ref.~\cite{huang2023bridging}, except that we simulate with open boundaries rather than periodic boundaries.


\section{Divergence of the detachment forces}

\subsection{Tissue fracture timescale depends systematically on time step}
\label{sec:ch4:timestep}

When performing a simple time-step convergence check of the AFV model, we find a strange phenomenon: when the tension difference $\Lambda > 0$ and a smaller time step $\Delta t$ is used to evolve the dynamics, cells appear to be more adherent to each other, and fewer gaps form between cells.
As shown in Fig.~\ref{fig:ch4:timestep}a, we evolve a cluster of $N=100$ cells using different time steps, $\Delta t=0.01$ and $\Delta t=0.005$.
After a simulation of $100$ time units, the cluster evolved with the smaller $\Delta t$ is more compact than the one evolved with the larger $\Delta t$, and unlike the large-$\Delta t$ cluster, it has not fractured or formed holes.
We have also confirmed that this behavior occurs in the code released by Ref.~\cite{huang2023bridging}. One way to quantify cluster cohesion is to consider the time scale of rupture of the cluster---i.e., for an initially $N$-cell cluster, due to the random motility of individual cells, when does it rupture into two disconnected clusters?
This can be measured using the survival probability $S(t)=\mathbb{P}(T>t)$, where we monitor the time $T$ of the \emph{first} rupture event~\cite{wang2026controlling, wang2025confinement}.
Fig.~\ref{fig:ch4:timestep}b shows $S(t)$ for a range of $\Delta t$: as $\Delta t$ decreases, rupture is delayed, and the survival curves shift to longer times.
A convenient summary statistic is the median survival time $t_{1/2}=S^{-1}(1/2)$---the time at which the survival probability drops to one half---which grows rapidly as $\Delta t$ is reduced (blue curve in Fig.~\ref{fig:ch4:timestep}c).

\begin{figure}[tp]
\centering
\includegraphics[width=0.85\textwidth]{chap4/timestep.pdf}
\caption{Tissue cohesion is strongly dependent on simulation timestep.
\textbf{a}, Simulation snapshots of an $N=100$ cell cluster at $t=100$ for time steps $\Delta t = 0.01$ (upper) and $\Delta t = 0.005$ (lower).
Each simulation was first evolved for $20$ time units with zero motility to reach a steady state, followed by $100$ time units of active dynamics.
\textbf{b}, Survival probability $S(t)$ of an initially intact monolayer with time steps $\Delta t=0.1$, $0.05$, $0.02$, $0.01$, $0.009$, $0.008$, $0.007$, $0.006$, and $0.005$.
\textbf{c}, Median survival time $t_{1/2}=S^{-1}(1/2)$ as a function of $\Delta t$; the cutoff is set to $\delta=0.45$ for the truncated curve.
For each value of $\Delta t$, we apply standard Kaplan-Meier survival analysis~\cite{kaplan1958nonparametric} to the results of $480$ independent simulations; error bars represent $95\%$ confidence intervals.
Parameters: number of cells $N=100$, initial packing fraction $\phi=0.5$, $\ell=1$, $K_P=1$, $A_0=\pi$, $P_0=4.8$, $\Lambda=0.1$, angular diffusion coefficient $D_r=1.33$, and active velocity $v_0=1.5$.}
\label{fig:ch4:timestep}
\end{figure}

Fig.~\ref{fig:ch4:timestep}b,c shows that our qualitative visual analysis is correct: the cells become more adherent as the time step $\Delta t$ decreases. Why does this happen? Does this reflect an issue with numerical algorithms, such as an incorrect choice of time integration schemes, or is it something more fundamental? In the next section, we show that this is an unavoidable consequence of the underlying geometric assumptions of the model.


\subsection{Forces between a cell doublet}
\label{sec:ch4:doublet}

To identify the origin of the divergence, we begin with the simplest detachment event: two identical cells that share a single contact.
As shown in Fig.~\ref{fig:ch4:doublet}a, we place the two cell centers symmetrically at $\mathbf{r}_\pm=\pm(\ell-\epsilon)\hat{x}$, so that the center-to-center distance is $d=2(\ell-\epsilon)$~\cite{teomy2018confluent}.
The parameter $\epsilon$ therefore measures how close the pair is to detachment: $\epsilon>0$ corresponds to a finite contact ($d<2\ell$), and $\epsilon\to0^+$ is the limit in which cells detach.
In this geometry, each cell boundary consists of a circular arc of length $P^{(n)}=\ell\phi$ together with a straight contact segment of length $P^{(c)}=2\sqrt{\ell^2-(\ell-\epsilon)^2}$, where $\phi=2\pi-2\operatorname{atan2}\bigl(\sqrt{\ell^2-(\ell-\epsilon)^2},\ell-\epsilon\bigr)$ is the angle spanning the region not in contact.

\begin{figure}[t]
\centering
\includegraphics[width=\textwidth]{chap4/doublet.pdf}
\caption{Cell doublet.
\textbf{a}, Schematic of a cell doublet with centers positioned symmetrically at $\mathbf{r}_\pm=\pm(\ell - \epsilon)\hat{x}$. The arc radius of each cell is $\ell$, and the distance between the two cell centers is given by $d = 2(\ell - \epsilon)$, thus the contact length (red) is $2\sqrt{\ell^2-(\ell-\epsilon)^2}$.
\textbf{b}, Force between a cell doublet as a function of their separation distance for $\Lambda=0,$ $0.1$, $0.2$, $0.3$, $0.4$, and $0.5$ (blue to red). Positive values $f>0$ correspond to attractive forces. Solid lines show Eq.~\eqref{eq:ch4:force}, and dashed lines are the asymptotic form for $d\to2\ell$ from Eq.~\eqref{eq:ch4:force-asymp}. Default parameters: $\ell=1$, $K_P=1$, $A_0=\pi$, and $P_0=4.8$.}
\label{fig:ch4:doublet}
\end{figure}

Because the doublet is symmetric, it is convenient to use the energy of a single cell:
\begin{equation}
\label{eq:ch4:single-cell}
    E_\mathrm{s}=K_A(A-A_0)^2+K_P(P-P_0)^2+\Lambda P^{(n)},
\end{equation}
where the area and perimeter of each cell are
\begin{eqnarray}
    A&=&(\ell-\epsilon)\sqrt{\ell^2-(\ell-\epsilon)^2}+\frac{\ell^2\phi}{2},\nonumber\\
    P&=&2\sqrt{\ell^2-(\ell-\epsilon)^2}+\ell\phi.\nonumber
\end{eqnarray}
The (passive) interaction force between the two cells follows from the derivative of $E_\mathrm{s}$ with respect to the separation parameter,
$f=-\partial E_\mathrm{s}/\partial\epsilon$.
Evaluating this derivative gives (see Sec.~\ref{sec:ch4:appendices} for details and some asymptotic limits)
\begin{equation}
\label{eq:ch4:force}
    f=4\sqrt{(2\ell-\epsilon)\epsilon}\left[K_A(A-A_0)+ K_P\dfrac{(P-P_0)}{2\ell-\epsilon}\right] + {\Lambda}\dfrac{2\ell}{\sqrt{(2\ell-\epsilon)\epsilon}},~~~(\ell\geqslant\epsilon>0)
\end{equation}
where in our sign convention $f>0$ corresponds to attractive forces (i.e., forces that tend to increase $\epsilon$). We see that the last term in Eq.~\eqref{eq:ch4:force} diverges as the two cells separate, $\epsilon \to 0^+$. This term comes from the diverging derivative of the non-contacting perimeter $P^{(n)}$ with $\epsilon$---in the finite Voronoi representation, small changes of the cell centers can lead to very large changes of perimeters.

The divergence of the detachment forces explains the time-step dependence of the rupture time scale: when $\Lambda > 0$, cells \emph{should not be able to detach from one another} in the FV model of Refs.~\cite{huang2023bridging, teomy2018confluent}. The presence of detachment occurs only because we are using a finite time step $\Delta t$, where the numerical evolution can allow the cell to ``skip'' the divergence to reach full cell-cell separation, where the force is zero.
It is possible for numerical errors to skip the divergence in part because the rapid increase in force is only apparent at distances very close to full separation (Fig.~\ref{fig:ch4:doublet}b)---if the cells move from distance $d\approx 1.95$ to over $2$ in a single time step, the divergence would not be felt. Given the value $v_0 = 1.5$ and a time step $\Delta t = 0.05$, this could be quite common.
By decreasing the time step $\Delta t$, we are approaching the ``correct'' no-detachment limit where $t_{1/2}\to\infty$ (Fig.~\ref{fig:ch4:timestep}c). Clearly, this behavior is not desired in a model of cells that can detach from one another, and we attempt to regularize these divergences in the next section.


\subsection{Regularization by introducing a cutoff}
\label{sec:ch4:regularization}

To regularize the divergent detachment force, it is useful to identify precisely where the divergence enters the force calculation in the active finite Voronoi model.
Given the cell-center positions of $N$ cells $\{\mathbf{r}_i\}$, the finite Voronoi structure produces a set of vertices $\{\mathbf{h}_m\}$ (including inner triple-junction vertices and outer vertices where straight edges meet circular arcs; see Fig.~\ref{fig:ch4:FV_model}).
Forces on cell centers follow from the chain rule:
\begin{equation}
\label{eq:ch4:chainrule}
    {f}_{i,x}=-\frac{\partial E}{\partial x_i}=-\sum_m\frac{\partial E}{\partial \mathbf{h}_m}\cdot\frac{\partial \mathbf{h}_m}{\partial x_i},
\end{equation}
where $x_i$ is the $x$ component of $\mathbf{r}_i$, with an analogous expression for the $y$ component.

The divergence does \emph{not} originate from the inner vertices (circumcenters of Delaunay triangles), whose derivatives remain finite for generic configurations.
Instead, it comes from \emph{outer} vertices $\mathbf{h}^{\mathrm{out}}$ that connect a straight Voronoi segment to a circular arc.
For two cells $i$ and $j$ sharing an outer vertex, there are two points $\mathbf{h}_\pm^\mathrm{out}$ that lie at a distance $\ell$ from both centers, which are~\cite{teomy2018confluent}
\begin{equation}
\label{eq:ch4:hout}
    \mathbf{h}^{\textrm{out}}_\pm=\frac{\mathbf{r}_i+\mathbf{r}_j}{2}\pm\frac{\sqrt{4\ell^2-|\mathbf{r}_i-\mathbf{r}_j|^2}}{2|\mathbf{r}_i-\mathbf{r}_j|}(\mathbf{r}_i-\mathbf{r}_j)\times\hat{z}.
\end{equation}
The $\pm$ reflects the two symmetric intersection points of the circles of radius $\ell$ centered at $\mathbf{r}_i$ and $\mathbf{r}_j$, lying on opposite sides of the line $ij$.
Their derivatives with respect to $x_i$ are given by
\begin{equation}
\label{eq:ch4:outer-deriv}
    \frac{\partial \mathbf{h}^\textrm{out}_\pm}{\partial x_i}=\dfrac{\hat{x}}{2}\mp\frac{2\ell^2(x_i-x_j)}{\sqrt{4\ell^2-|\mathbf{r}_i-\mathbf{r}_j|^2}}\frac{(\mathbf{r}_i-\mathbf{r}_j)}{|\mathbf{r}_i-\mathbf{r}_j|^3}\times\hat{z}\mp\frac{\sqrt{4\ell^2-|\mathbf{r}_i-\mathbf{r}_j|^2}}{2|\mathbf{r}_i-\mathbf{r}_j|}\hat{y}.
\end{equation}
We can see that $\partial \mathbf{h}_\pm^{\mathrm{out}}/\partial x_i$ contains a term with $\sqrt{4\ell^2-|\mathbf{r}_i-\mathbf{r}_j|^2}$ in the denominator, which diverges as $|\mathbf{r}_i-\mathbf{r}_j|\to2\ell$.
This geometric divergence is the origin of the divergent detachment force in Eq.~\eqref{eq:ch4:force}.

Physically, $\sqrt{4\ell^2-|\mathbf{r}_i-\mathbf{r}_j|^2}=2\sqrt{\ell^2-(|\mathbf{r}_i-\mathbf{r}_j|/2)^2}$ is exactly the distance between the two outer vertices $\mathbf{h}^\textrm{out}_{\pm}$ generated by the pair $\{\mathbf{r}_i,\mathbf{r}_j\}$, as defined in Eq.~\eqref{eq:ch4:hout}.
For the cell doublet shown in Fig.~\ref{fig:ch4:doublet}a, this quantity is simply the contact length $P^{(c)}$ (red).
Thus, a minimal practical regularization is to bound the vanishing denominator in Eq.~\eqref{eq:ch4:outer-deriv} by a small cutoff $\delta$:
\begin{equation}
\label{eq:ch4:cutoff}
    \sqrt{4\ell^2-|\mathbf{r}_i-\mathbf{r}_j|^2} \to \max\left(\sqrt{4\ell^2-|\mathbf{r}_i-\mathbf{r}_j|^2},\delta\right).
\end{equation}
This prescription renders forces finite at detachment while leaving them unchanged for $|\mathbf{r}_i-\mathbf{r}_j|<\sqrt{4\ell^2-\delta^2}=2\ell-(\delta^2/4\ell)+\mathcal{O}(\delta^4)$, at the cost of introducing a new microscopic length scale $\delta$. This cutoff can be interpreted as a threshold contact length below which the Voronoi shape assumption is faulty, analogous to the finite edge-length criterion used to trigger T1 transitions in vertex models~\cite{chen2022activation, bi2015density, fletcher2014vertex, spahn2013vertex}.

\begin{figure}[ht]
\centering
\vspace{0.3cm}
\includegraphics[width=0.75\textwidth]{chap4/truncated.pdf}
\caption{Regularization by introducing a cutoff.
\textbf{a}, Force between a cell doublet as a function of cell-center distance for $\Lambda=0.2$ and $P_0=4.8$. Empty circles and triangles represent forces obtained from simulations before and after implementing the truncation; analytical results are the corresponding curves in Fig.~\ref{fig:ch4:doublet}b. Inset: enlarged view around the cutoff $\delta=0.45$, which corresponds to $d_c=\sqrt{4\ell^2-\delta^2}\approx 1.95$.
\textbf{b}, Phase diagram for the final detachment force $f_\mathrm{detach}$ when $d=2\ell$ in the $P_0$--$\Lambda$ plane, with cutoff $\delta=0.45$.
Default parameters: $\ell=1$, $K_P=1$, and $A_0=\pi$.}
\label{fig:ch4:truncated}
\end{figure}

From the shapes of the curves in Fig.~\ref{fig:ch4:doublet}b, for $\Lambda>0$ the force-distance relation exhibits a local minimum near $d\approx2\ell$, after which the force increases rapidly toward divergence.
For typical parameters, this rapid growth sets in around $d=1.95\ell$, corresponding to a threshold value $\delta=\sqrt{4\ell^2-d^2}\approx 0.45\ell$.
Applying the cutoff of Eq.~\eqref{eq:ch4:cutoff} with $\delta = 0.45$ (as $\ell = 1$), we successfully regularize the detachment forces around $d=2\ell$ (see Fig.~\ref{fig:ch4:truncated}a), and the truncation effect kicks in at $d_c\approx1.95$ as expected (inset of Fig.~\ref{fig:ch4:truncated}a). With this truncation, we eliminate the time-step dependence of the tissue fracture timescale (Fig.~\ref{fig:ch4:timestep}c, red curve).

In the finite Voronoi model, cell doublets are forced to detach when the cell-center distance $d$ exceeds $2\ell$. We extract the final detachment forces $f_\textrm{detach}$ at $d=2\ell$ numerically, and we find that, after introducing the regularization, this force is no longer divergent even when sweeping through the parameter space of $P_0$ and $\Lambda$ (Fig.~\ref{fig:ch4:truncated}b). The detachment forces in Fig.~\ref{fig:ch4:truncated}b increase strongly with $\Lambda$, as we would expect---increasing $\Lambda$ corresponds to a higher energy cost for the cell-medium interface, increasing the cost for separation. We also see that increasing $P_0$ while holding other parameters constant weakly decreases the detachment force. This is consistent with, e.g., looking at the analytical detachment force in Eq.~\eqref{eq:ch4:force} at the distance at which regularization kicks in, $\epsilon_c = \ell - d_c/2$.


\section{Calibration with a deformable polygon model}
\label{sec:ch4:calibration}

The origin of the divergence we found in the finite Voronoi model is the assumption that, even as the cells are almost completely separated, the cells remain circular with a flat contact.
During detachment, however, real cells exhibit appreciable shape deformation~\cite{smeets2019effect, vangheel2026rigidity, chu2005johnson, byers1995role, rozema2025remodeling,eckert2025cell}.
While we are able to regularize the divergence by introducing a new scale $\delta$, we want to better understand how to set this value. To calibrate our regularized finite Voronoi model, we compare it against a deformable polygon (DP) model that has the same energy function but in which cell boundaries are explicitly represented by polygons with large numbers of vertices that can deform continuously~\cite{boromand2018jamming, lv2024active, weng2022convergent}---i.e., the same energetics as our FV model but without the Voronoi geometry assumption. During calibration, we focus on passive mechanical forces and set $v_0=0$.

\subsection{Matching passive steady states between the deformable polygon model and finite Voronoi model}
\label{sec:ch4:DP-model}

In the DP model, each cell is represented by a closed polygon with an ordered set of vertices $\{\mathbf h_{m}\}_{m=1}^M$.
We use the same energy as in the FV model: quadratic area and perimeter elasticity together with distinct line tensions on cell-cell contacts and cell-medium boundary.
However, the degrees of freedom are not the cell centers, but the vertices on the cell boundaries.
Vertex dynamics are evolved by overdamped relaxation:
\begin{equation}
\label{eq:ch4:dp-evolve}
    \frac{\mathrm{d}\mathbf{h}_m}{\mathrm{d}t}=-\mu\frac{\partial E}{\partial\mathbf{h}_m}.
\end{equation}

To place the two models on a comparable footing, we first choose parameters so that the steady states of passive cell doublets match between the models. Minimizing the energy in the DP model using Eq.~\eqref{eq:ch4:dp-evolve}, we see that the minimal-energy state is, as the finite Voronoi model assumes, generally a pair of circular arcs with a straight boundary between them. However, the resulting steady-state radius of the arcs is the value that minimizes the energy, which will depend on the mechanical parameters of the model---and not the default $\ell=1$ used in the FV model. To match the FV and DP models, we must vary $\ell$ when we vary the FV model parameters, choosing the value that minimizes the energy. In this case, as long as the steady state of the DP model fits the finite Voronoi assumption, the FV model should recover the same steady state. Therefore, in our calibration process, we find the minimal-energy separation between cell centers and the minimal-energy value of $\ell$,
\begin{equation}
\label{eq:ch4:optimal}
(\ell_0,\epsilon_0)=\operatorname*{arg\,min}_{\ell\geqslant\epsilon>0} E_\mathrm{s}(\ell,\epsilon),
\end{equation}
and then set the FV radius to $\ell=\ell_0$. We see that in the absence of external forcing, FV and DP predict the same relaxed doublet shape (Fig.~\ref{fig:ch4:steady_state}a).
The values of the steady-state radius $\ell_0$ for different sets of $(P_0,\Lambda)$ are shown in Fig.~\ref{fig:ch4:steady_state}b. We see that $\ell_0$ increases with the preferred perimeter $P_0$, while $\Lambda$ has a relatively weaker effect on $\ell_0$.

\begin{figure}[t]
\centering
\includegraphics[width=0.85\textwidth]{chap4/steady_state.pdf}
\caption{Steady states of deformable polygon model.
\textbf{a}, Simulation snapshots of how the cell doublet relaxes from its initial shape to the steady state in DP model. Blue shapes indicate the optimal $(\ell_0, \epsilon_0)\approx(0.87, 0.12)$ computed from Eq.~\eqref{eq:ch4:optimal} for $P_0=4.8$ and $\Lambda=0.2$.
\textbf{b}, Phase diagram for optimal $\ell_0$ in the $P_0$--$\Lambda$ plane.
Default parameters: $K_P=1$, $A_0=\pi$.}
\label{fig:ch4:steady_state}
\end{figure}


\subsection{Force-distance curves in the deformable polygon model}

We want to find the relationship of force between cell pairs and distance in the DP model, akin to the values predicted for the finite Voronoi model in Fig.~\ref{fig:ch4:doublet}b. To do this, we apply equal and opposite forces to each cell in an initially relaxed doublet in the DP model, and measure the resulting displacement of the cell centroids $d$. (Note that to be consistent in our sense of cell-cell separation between the DP and FV model, here and throughout this section, we use $d$ to denote the centroid-centroid separation of both the DP and FV cells---cells in the DP model are not perfectly circular and do not have a ``center'' defined as in the FV model.)
We distribute the force on one cell uniformly over all vertices of each cell to avoid spurious torques~\cite{smeets2019effect}, as shown in Fig.~\ref{fig:ch4:detachment}a.
We then record the centroid separation and contact length $P^{(c)}$ (red segments in Fig.~\ref{fig:ch4:doublet}a) as functions of the external force $f$. We define the detachment force $f_{\textrm{detach}}^\textrm{DP}$ as the force at which the contact length falls below the typical distance $\ell_c=2\pi\ell_0/M$ between neighboring vertices (see Sec.~\ref{sec:ch4:appendices}).
The detachment force increases strongly with $\Lambda$ as expected, and decreases weakly with $P_0$ (Fig.~\ref{fig:ch4:detachment}b).
If we choose parameters of the FV model such that the steady states of the FV and DP models match (i.e., $\ell$ is the energy-minimizing $\ell_0$), we find that the FV detachment force exhibits the opposite trend with $P_0$: for fixed $\Lambda$, increasing $P_0$ leads to larger $f_\textrm{detach}^\textrm{FV}$ (Fig.~\ref{fig:ch4:detachment}c), in contrast to the behavior in the DP model (Fig.~\ref{fig:ch4:detachment}b) and in the FV model with a fixed maximum radius ($\ell=1$, Fig.~\ref{fig:ch4:truncated}b).

What behavior should we expect for detachment forces' dependence on $P_0$? In our rescaled units, $P_0=-\lambda^{(c)}/2K_P$, so increasing $P_0$ typically corresponds to stronger cell-cell adhesion (i.e., making cell-cell line tension $\lambda^{(c)}$ more negative). Stronger cell-cell adhesion must increase the detachment force, so at first glance, it initially appears that only Fig.~\ref{fig:ch4:detachment}c exhibits the physically correct trend.
However, $P_0$ and $\Lambda=\lambda^{(n)}-\lambda^{(c)}$ both depend on the more fundamental parameter $\lambda^{(c)}$.
If $\lambda^{(n)}$ is held fixed and only $\lambda^{(c)}$ is varied, increasing adhesion corresponds to moving along the line $\Lambda=\lambda^{(n)}+2K_P P_0$ in the $\Lambda$--$P_0$ plane.
Along this path, the detachment forces in both the finite Voronoi model and the deformable polygon model increase as cell-cell adhesion is made stronger (Fig.~\ref{fig:ch4:detachment}d).

We compute the force-distance curve in the DP model and track the contact length as the cells are pulled apart from one another. We show these curves in Fig.~\ref{fig:ch4:calibrate}a,b for two different values of $P_0$. Here $f>0$ corresponds to a pull-off force dipole while $f<0$ corresponds to a pushing force dipole.
The contact length $P^{(c)}$ of the cell doublet decreases roughly linearly with the cell-centroid distance, and in the pulling limit ($f>0$) the force for the DP model increases with centroid distance almost linearly, consistent with Ref.~\cite{smeets2019effect}, where cells are modeled as three-dimensional elastic triangulated shells.

Comparing the analytical curves of the FV model and the DP simulation results in Fig.~\ref{fig:ch4:calibrate}a,b, we find two key results: 1) for small perturbations away from the relaxed state (zero force), the DP and FV models agree well, and 2) the DP model does not have a divergent detachment force, staying attached over a much larger cell-cell distance than the FV model. We also see by looking at the cell doublet shapes in Fig.~\ref{fig:ch4:detachment}a that at larger forces, the finite Voronoi shape assumption starts to fail---cells can have centroids separated by more than $2\ell_0$ without losing contact, and at higher forces cells are no longer well described by circular arcs (rightmost panel of Fig.~\ref{fig:ch4:detachment}a).

We have found that, simply by choosing the energy-minimizing $\ell_0$, the small forces near the equilibrium are in agreement between FV and DP models. However, the detachment forces are clearly not in good agreement. Fig.~\ref{fig:ch4:calibrate} shows two potential approaches to address this.

\begin{figure*}[tp]
\centering
\includegraphics[width=0.9\textwidth]{chap4/detachment.pdf}
\caption{Detachment of cell doublets in the finite Voronoi model (FV) and deformable polygon (DP) models.
\textbf{a}, Snapshots of a cell doublet for both deformable polygon model and finite Voronoi model at different centroid distances.
Preferred perimeter is set to $P_0=4.8$ and $\Lambda=0.2$, yielding the optimal $(\ell_0, \epsilon_0)\approx (0.87, 0.12)$ which has been used to set the maximum radius $\ell=\ell_0$ in FV model. The four centroid distances from left to right are $1.374$ (compressed), $1.550$ (steady state), 1.735 (approaching detachment threshold $2\ell_0$ in FV), and $1.912$ ($>2\ell_0$).
\textbf{b}, The detachment forces $f_\textrm{detach}^\textrm{DP}$ for the DP model where contact length $P^{(c)}$ of the cell doublet approaches zero.
\textbf{c}, The detachment forces $f_\textrm{detach}^\textrm{FV}$ for the FV model when $d=2\ell$, where we have set the maximum radius to the steady-state radius from the DP model, i.e., $\ell=\ell_0$, and the threshold value $\delta=0.45$.
\textbf{d}, Detachment forces along the path $\Lambda=\lambda^{(n)}+2K_P P_0$ in the previous phase diagrams, with the line tension on cell-medium interfaces fixed at $\lambda^{(n)}=-7.9$.
Default parameters: $K_P=1$, $A_0=\pi$.}
\label{fig:ch4:detachment}
\end{figure*}


\subsection{Calibration strategy 1: setting a critical length}
\label{sec:ch4:cal1}

As a first calibration strategy for near-detachment behavior, we choose the energy-minimizing value $\ell = \ell_0$ and adopt a fixed value of $\delta$, using the default value $\delta=0.45$ as in Fig.~\ref{fig:ch4:truncated}. Choosing a fixed value of $\delta$ corresponds to the idea that there is a minimal contact length beyond which cell-cell forces cannot increase, or equivalently a maximal cell-cell separation.
We choose the value $\delta = 0.45$ to capture a reasonable scale at which the force-distance curve flattens as in Fig.~\ref{fig:ch4:truncated}a---in strategy 1 we do not use any information from the DP model to set this value.
After setting $\ell=\ell_0$ and applying this fixed cutoff to truncate the rapidly growing force near detachment $d\approx2\ell_0$ in the FV model,
the simulation results (empty squares in Fig.~\ref{fig:ch4:calibrate}a,b) follow the analytical curves but lack the divergence near $2\ell_0$, as expected. Therefore, consistent with the comparison in the previous section between the analytical FV results and the DP model, we observe reasonable agreement between the FV and DP force-distance relations for $d\leqslant 2\ell_0$, with only a few deviations at the smallest cell-cell distances where cells are highly compressed. However, the agreement in the force-distance curves necessarily means that the detachment forces differ between DP and FV in strategy 1---since cells remain connected even when $d>2\ell_0$ in the DP model, the DP model exhibits a much larger detachment force.

Given this strategy, how does cluster rupture, as we explored in Fig.~\ref{fig:ch4:timestep}b,c, change with mechanical parameters like $P_0$, $v_0$, and $\Lambda$? We study the full active finite Voronoi model with this truncation and calibration procedure. As shown in Fig.~\ref{fig:ch4:calibrate}c, for fixed $v_0$ and $P_0$, increasing $\Lambda$ consistently makes clusters harder to fracture. This is because increasing $\Lambda$ effectively increases $\lambda^{(n)}$, thereby penalizing the creation of non-contacting free boundaries.
In standard active Voronoi and vertex models, it is well established that larger preferred perimeter or higher motility $v_0$ leads to more fluid-like behavior.
Our result shows that holding all other parameters constant, a larger preferred perimeter $P_0$---corresponding to a more fluid-like state with a lower line tension $\lambda^{(c)}$ for contacting edges---makes clusters harder to rupture (Fig.~\ref{fig:ch4:calibrate}d). This differs from Ref.~\cite{huang2023bridging}, where clusters with larger $P_0$ were more likely to be in dispersed states.
We attribute this discrepancy not to the choice of cutoff $\delta$, but to calibrating the maximum cell radius $\ell$ to its energy-minimizing value $\ell_0$, which leads to detachment forces that increase with $P_0$ (Fig.~\ref{fig:ch4:detachment}c).


\subsection{Calibration strategy 2: matching detachment forces}
\label{sec:ch4:cal2}

As discussed above, the previous strategy aligns the FV and DP models only for $d\leqslant 2\ell_0$, and this results in the FV model requiring a much smaller pull-off force to detach two cells.
Thus, we also consider an alternative calibration based on the detachment force $f_\textrm{detach}$.
In this approach, we set $\ell$ to be the energy-minimizing $\ell_0$ and adjust $\delta$ such that the detachment force in the FV model matches that of the DP model (Fig.~\ref{fig:ch4:detachment}b), i.e., $f_\textrm{detach}^\textrm{FV}=f_\textrm{detach}^\textrm{DP}$.
We numerically compute the value of $d_c$ required for Eq.~\eqref{eq:ch4:force} to equal $f_\textrm{detach}^\textrm{DP}$, and obtain the corresponding $\delta_c$. The results are shown in Fig.~\ref{fig:ch4:calibrate}e,f, where the horizontal dashed lines indicate the matched detachment forces.

As shown in Fig.~\ref{fig:ch4:calibrate}e,f, this calibration strategy introduces sharply increasing forces near detachment which will naturally require finer timestepping; to obtain reliable fracture statistics, we therefore gradually reduce the time step $\Delta t$ from the default value of $0.01$ to $0.002$ until the results converge.
In the phase diagram of the median survival time $t_{1/2}$ (Fig.~\ref{fig:ch4:calibrate}g,h), this calibration generally makes clusters harder to detach than under the previous strategy (more points appear as stars), since the detachment forces $f_\textrm{detach}^\textrm{FV}$ are much larger now.
Moreover, because the detachment forces are matched to those of the DP model, the transition lines between no-rupture and rupture states differ from those in Fig.~\ref{fig:ch4:calibrate}c,d: the fracture timescale depends only weakly on the preferred perimeter $P_0$ now, consistent with the detachment-force phase diagram in Fig.~\ref{fig:ch4:detachment}b, whereas the previous strategy (Fig.~\ref{fig:ch4:calibrate}c,d) corresponds to Fig.~\ref{fig:ch4:detachment}c.
The other trends remain similar to the earlier results: decreasing $\Lambda$ or increasing $v_0$ both promote fracture.

\begin{figure*}[t]
\centering
\includegraphics[width=\textwidth]{chap4/calibrate.pdf}
\caption{Calibrating the finite Voronoi model (FV) to a deformable polygon (DP) model.
\textbf{a},\textbf{b} show how the forces in both models and also the contact length $P^{(c)}$ in DP model (green) vary with increasing centroid distance at two different $P_0$ values ($\Lambda=0.2$ and $\delta=0.45$).
\textbf{c},\textbf{d} are the phase diagrams of median survival time $t_{1/2}$ on $P_0$--$\Lambda$ and $P_0$--$v_0$ planes, where red crosses ($t_{1/2}=0$) indicate tissues are already ruptured at the start (after initial relaxation time), and purple stars indicate that $t_{1/2}$ exceeds the total simulation time. Threshold value is set to $\delta=0.45$.
\textbf{e}--\textbf{h}, Calibrate FV model by matching the detachment forces $f_\textrm{detach}$ to DP model with varying cutoff $\delta$.
Gray dashed lines in \textbf{e},\textbf{f} indicate the match of detachment forces in both models.
Each data point of $t_{1/2}$ in the phase diagrams is obtained by computing the survival probability from the results of $480$ independent simulations.
Default parameters: $K_P=1$, $A_0=\pi$, $D_r=1.33$.}
\label{fig:ch4:calibrate}
\end{figure*}


\section{Discussion}
\label{sec:ch4:discussion}

Our results show that the finite Voronoi model contains a pathological behavior when cell-cell and cell-medium tensions differ: the force required to separate two cells diverges, so fracture in active clusters becomes spuriously controlled by the numerical time step rather than by physical parameters.
We trace the origin of this divergence analytically by analyzing a two-cell system and introduce a cutoff threshold that removes the divergence and restores a well-defined fracture timescale. We then compare the finite Voronoi model with a deformable polygon model and propose two calibration strategies, providing a practical framework for regularizing and calibrating the model in studies of tissue fracture and cell separation.

The cutoff regularization introduced in Eq.~\eqref{eq:ch4:cutoff} removes the time step dependence of the AFV model, but it modifies near-detachment mechanics by introducing a new length scale $\delta$.
We think of $\delta$ as a truncation threshold for the contact length $P^{(c)}$: once the contact shrinks below $\delta$, the intercellular force no longer increases.
However, the precise value of this threshold is both model- and parameter-dependent, and cannot be uniquely fixed within the FV framework.
Recent work using finite-element descriptions of cell mechanics has emphasized that tissue-level phenomena such as jamming can depend sensitively on the force required to separate neighboring cells~\cite{vangheel2026rigidity}, suggesting that the detachment force should be viewed as an important physical control parameter rather than a purely numerical detail.
From this perspective, the regularization parameter $\delta$ provides a practical way to tune the FV model so that its near-detachment mechanics match those of more detailed descriptions.
For the two calibration strategies introduced in Sec.~\ref{sec:ch4:calibration}, the appropriate choice depends on the physical regime of interest.
If detachment events are rare---for example, when the tissue remains largely in a solid-like state---strategy~1 is typically sufficient, since it already reproduces the force-distance relationship well near the steady-state configuration.
By contrast, when frequent cell rearrangements or fracture events occur, strategy~2 is more appropriate because it matches the detachment forces between the FV and DP models.
More generally, the cutoff $\delta$ can also be calibrated to match detachment forces obtained from more detailed models, such as those in Refs.~\cite{boromand2018jamming, lv2024active}, or even from experimental measurements of cell-cell pull-off forces.

The divergence of detachment forces in the finite Voronoi model may affect the results of Ref.~\cite{huang2023bridging} for simulations done with $\Lambda > 0$.
First, because no force cutoff was implemented, clusters can fracture only due to the finite simulation time step $\Delta t$. Although the numerical dynamics with finite $\Delta t$ will act as an effective regulator similar to the threshold $\delta$ introduced here, we think the use of the regularizing $\delta$ avoids possible complications in which varying numerical details like $\Delta t$ can have spurious physical effects. We also expect some quantitative differences between regularization with a finite $\Delta t$ and our approach.
Secondly, choices about the calibration of $\ell$ and regularization can qualitatively change results. Using a fixed $\ell$ rather than the steady-state radius $\ell_0$ from the DP model reverses the trend of fracture behavior when varying $P_0$ (Figs.~\ref{fig:ch4:truncated}b and \ref{fig:ch4:detachment}c).
Interestingly, in our model using a varying $\ell_0$ and calibrating to the detachment forces (Fig.~\ref{fig:ch4:calibrate}h), which we think is the most reasonable approach, we do see the same qualitative trends in the effect of $P_0$ that are observed by Ref.~\cite{huang2023bridging} in their phase diagrams (e.g., their Fig.~2b). However, we believe the origin of these behaviors may be subtler than expected, and they depend on our calibration choices.

The associated purpose-built Python package \texttt{PyAFV} is available on GitHub at \href{https://github.com/wwang721/pyafv}{wwang721/pyafv}. A snapshot of the package and other data and code required to reproduce this chapter have been archived on Zenodo, doi:~\href{https://doi.org/10.5281/zenodo.18091659}{10.5281/zenodo.18091659}.


\section{Appendices}
\label{sec:ch4:appendices}

\subsection{Simulation details}

We simulate the dynamics of both the FV and DP models using the forward Euler (Euler-Maruyama) method~\cite{kloeden1992stochastic}.
Default parameters are given in Table~\ref{tab:ch4:params}. All parameters are nondimensionalized using the length scale $L=\sqrt{A_0/\pi}$, energy scale $K_A L^4$, and time scale $1/\mu K_A L^2$.
Note that we are not using the maximum radius $\ell$ as the scaling length unit so that $\ell$ can be set to other values than $1$, such as the optimal $\ell_0$ from minimizing $E_s$ from Eq.~\eqref{eq:ch4:single-cell}.

When simulating the dynamics of the finite Voronoi model, before activating the active motility $v_0$, the initial cell center positions are randomly distributed within a square domain at a packing fraction of about 0.5~\cite{huang2023bridging}. The system is then evolved for $T_\textrm{relax}=20$ time units to reach a steady state. Subsequently, active motility is turned on, and the simulation is run for a duration $T$.
The time since relaxation is used in the survival probability and mean survival time calculations.

\begin{table}[ht]
\centering
\caption{Default simulation parameters used throughout this chapter; any deviations are explicitly specified.}
\label{tab:ch4:params}
\begin{tabular}{llcc}
\toprule\toprule
Parameter & \quad Description & Scale & Value\\\midrule
$A_0$ & Preferred area & $L^2$ & $\pi$\\
$P_0$ & Preferred perimeter & $L$ & $4.8$\\
$\ell$ & Maximum radius & $L$ & $1$\\
$K_P$ & Perimeter elastic modulus & $K_A L^2$ & $1$\\
$\Lambda$ & Tension difference $\lambda^{(n)}-\lambda^{(c)}$ & $K_A L^3$ & $0.2$\\
$\Delta t$ & Time step & $1/\mu K_A L^2$ & $0.01$\\
$T_\textrm{relax}$ & Initial relaxation time& $1/\mu K_A L^2$ & $20$\\
$T$ & Total simulation time & $1/\mu K_A L^2$ & $1000$\\
$D_r$ & Rotational noise & $\mu K_A L^2$ & $1.33$\\
$N$ & Number of cells & 1 & $100$\\
\bottomrule\bottomrule
\end{tabular}
\end{table}

\subsubsection{Details of finite Voronoi model force calculations}

In the standard/finite Voronoi model, as given in Refs.~\cite{bi2016motility, teomy2018confluent}, an inner vertex connecting three cells $i$, $j$, and $k$ is given by
\begin{equation}
    \mathbf{h}^\textrm{in}=\alpha_i\mathbf{r}_i + \alpha_j\mathbf{r}_j + \alpha_k\mathbf{r}_k.
\end{equation}
The three barycentric coordinates of the circumcenter are
\begin{subequations}
\begin{eqnarray}
    \alpha_i=|\mathbf{r}_j-\mathbf{r}_k|^2(\mathbf{r}_i-\mathbf{r}_j)\cdot(\mathbf{r}_i-\mathbf{r}_k)/D,\\
    \alpha_j=|\mathbf{r}_i-\mathbf{r}_k|^2(\mathbf{r}_j-\mathbf{r}_i)\cdot(\mathbf{r}_j-\mathbf{r}_k)/D,\\
    \alpha_k=|\mathbf{r}_i-\mathbf{r}_j|^2(\mathbf{r}_k-\mathbf{r}_i)\cdot(\mathbf{r}_k-\mathbf{r}_j)/D,
\end{eqnarray}
\end{subequations}
where $D=2|(\mathbf{r}_i-\mathbf{r}_j)\times(\mathbf{r}_j-\mathbf{r}_k)|^2$. Thus, its derivative with respect to $x_i$ is given by
\begin{equation}
    \frac{\partial\mathbf{h}^\textrm{in}}{\partial x_i}=\alpha_i\hat{x}+\frac{\partial\alpha_i} {\partial x_i}\mathbf{r}_i+\frac{\partial\alpha_j}{\partial x_i}\mathbf{r}_j+\frac{\partial\alpha_k}{\partial x_i}\mathbf{r}_k =\alpha_i\hat{x}+\frac{\partial\alpha_j}{\partial x_i}(\mathbf{r}_j-\mathbf{r}_i)+\frac{\partial\alpha_k}{\partial x_i}(\mathbf{r}_k-\mathbf{r}_i),
\end{equation}
where in the second equation we have used the relation $\alpha_i + \alpha_j+ \alpha_k=1$.
We note this corrects a typo in Ref.~\cite{teomy2018confluent}.
A similar formula naturally holds for the derivative with respect to $y_i$.
Then we need to compute
\begin{equation}
    \frac{\partial \alpha_j}{\partial \mathbf{r}_i} = \alpha_j\left[\frac{\mathbf{r}_k-\mathbf{r}_j}{(\mathbf{r}_j-\mathbf{r}_i)\cdot(\mathbf{r}_j-\mathbf{r}_k)} + 2\frac{\mathbf{r}_i-\mathbf{r}_k}{|\mathbf{r}_i-\mathbf{r}_k|^2}-2\frac{(\mathbf{r}_j-\mathbf{r}_k)\times\hat{z}}{[(\mathbf{r}_i-\mathbf{r}_j)\times(\mathbf{r}_j-\mathbf{r}_k)]_z}\right].
\end{equation}
The expression for $\partial \alpha_k/\partial \mathbf{r}_i$ follows by exchanging the indices $j$ and $k$.
We have also explicitly shown the expressions for outer vertices in Eqs.~\eqref{eq:ch4:hout} and \eqref{eq:ch4:outer-deriv}.

The mechanical force experienced by cell $i$ is given by the gradient of the energy:
\begin{equation}
    f_{i,x}=-\frac{\partial E}{\partial x_i}=-\sum_j\left[2K_A(A_j-A_0)\frac{\partial A_j}{\partial x_i}+2K_P(P_j-P_0)\frac{\partial P_j}{\partial x_i} + \Lambda\frac{\partial P_j^{(n)}}{\partial x_i}\right].
\end{equation}
We note that each cell can be decomposed into a polygon (connecting all vertices belonging to that cell), together with additional circular segments (a segment is defined as the circular sector minus the isosceles triangle formed by the two radii and the chord), i.e., $A_j=A_j^\textrm{poly}+A_j^\textrm{seg}$ and $P_j=P_j^\textrm{poly}+P_j^{(n)}$.
Thus, the derivatives of the area and perimeter can be decomposed into two corresponding parts:
\begin{equation}
    \frac{\partial A_j}{\partial x_i}=\frac{\partial A_j^\textrm{poly}}{\partial x_i}+\frac{\partial A_j^\textrm{seg}}{\partial x_i},\qquad\frac{\partial P_j}{\partial x_i}=\frac{\partial P_j^\textrm{poly}}{\partial x_i}+\frac{\partial P_j^{(n)}}{\partial x_i}.
\end{equation}
For the polygon part,
\begin{equation}
    \frac{\partial A_j^\textrm{poly}}{\partial x_i} = \sum_m\frac{\partial A_j^\textrm{poly}}{\partial \mathbf{h}_m}\cdot\frac{\partial \mathbf{h}_m}{\partial x_i},\qquad\frac{\partial P_j^\textrm{poly}}{\partial x_i} = \sum_m\frac{\partial P_j^\textrm{poly}}{\partial \mathbf{h}_m}\cdot\frac{\partial \mathbf{h}_m}{\partial x_i},
\end{equation}
where the sum runs over all vertices (both inner and outer), as in Eq.~\eqref{eq:ch4:chainrule}. The derivatives $\partial\mathbf{h}_m/\partial x_i$ are already given above, while the derivatives ${\partial A_j^\textrm{poly}}/{\partial \mathbf{h}_m}$ and ${\partial P_j^\textrm{poly}}/{\partial \mathbf{h}_m}$ are provided in Eq.~(B15) of Ref.~\cite{teomy2018confluent}.

We now compute the contribution from the segment part. For a cell $j$ with a total segment area $A_j^\textrm{seg}$ and total arc length $P_j^{(n)}$, consider a segment between two neighboring outer vertices $\mathbf{h}^\textrm{out}_m$ and $\mathbf{h}^\textrm{out}_{m+1}$ (ordered clockwise). Let $\phi_m^{j}$ denote the angle between the vector $\mathbf{u}=\mathbf{h}^\textrm{out}_m - \mathbf{r}_j$ and the $x$-axis, i.e., $\phi_m^j=\operatorname{atan2}(u_y, u_x)\in(-\pi,\pi]$. The arc length between the two vertices is then $\Delta P=\ell\Delta\phi$, and the area of the corresponding segment is
\begin{equation}
    \Delta A = \frac{\ell^2}{2}(\Delta\phi-\sin\Delta\phi),
\end{equation}
where $\Delta\phi\equiv \phi_m^{j}-\phi_{m+1}^{j}\pmod{2\pi}$. Thus we have
\begin{equation}
    \frac{\partial P_j^{(n)}}{\partial\phi_m^{j}}=\ell,\qquad\frac{\partial A^\textrm{seg}_j}{\partial\phi_{m}^{j}}=\frac{\ell^2}{2}(1-\cos\Delta\phi),
\end{equation}
and ${\partial P_j^{(n)}}/{\partial\phi_{m+1}^{j}} = -{\partial P_j^{(n)}}/{\partial\phi_{m}^{j}}$, ${\partial A^\textrm{seg}_j}/{\partial\phi_{m+1}^{j}}=-{\partial A^\textrm{seg}_j}/{\partial\phi_{m}^{j}}$.
Because an outer vertex such as $\mathbf{h}_m^\textrm{out}$ is shared by cell $j$ and a neighboring cell $i$ [see Eq.~\eqref{eq:ch4:hout}], the angle $\phi_m^{j}$ depends on both $\mathbf{r}_i$ and $\mathbf{r}_j$. With the definition $\mathbf{u}=\mathbf{h}_m^\textrm{out}-\mathbf{r}_j=|\mathbf{u}|(\cos\phi_m^j,\sin\phi_m^j)$ and $|\mathbf{u}|=\ell$, we have
\begin{equation}
    \frac{\partial \mathbf{u}}{\partial x_i}=\frac{\partial \mathbf{h}_m^\textrm{out}}{\partial x_i}=\ell(-\sin\phi_m^j,\cos\phi_m^j)\frac{\partial\phi_m^j}{\partial x_i},
\end{equation}
where $(-\sin\phi_m^j,\cos\phi_m^j)=-(\mathbf{u}\times\hat{z})/|\mathbf{u}|$ is a unit vector perpendicular to $\mathbf{u}$.
Thus, taking the dot product with $(-\sin\phi_m^j,\cos\phi_m^j)/\ell$, we obtain the derivative of $\phi_m^j$ with respect to $x_i$:
\begin{equation}
    \frac{\partial\phi_m^{j}}{\partial x_i}=-\frac{1}{\ell^2}[(\mathbf{h}_m^\textrm{out}-\mathbf{r}_j)\times\hat{z}]\cdot\frac{\partial\mathbf{h}_m^\textrm{out}}{\partial x_i},
\end{equation}
and the $y$-component has an analogous form.
Since $\mathbf{u}$ (and hence $\phi^j_m$) depends only on the relative position $\mathbf{r}_i-\mathbf{r}_j$, we also have $\partial{\phi^{j}_m}/\partial \mathbf{r}_j=-\partial{\phi^{j}_m}/\partial \mathbf{r}_i$. With all these derivatives in hand, we can then compute the segment contributions to the derivative of the area and perimeter:
\begin{eqnarray}
    \frac{\partial A_j^\textrm{seg}}{\partial x_i}=\sum_{\mathbf{h}_m^\textrm{out}}\frac{\partial A_j^\textrm{seg}}{\partial\phi_m^{j}}\frac{\partial\phi^{j}_m}{\partial x_i},\qquad\frac{\partial P_j^{(n)}}{\partial x_i}=\sum_{\mathbf{h}^\textrm{out}_m}\frac{\partial P_j^{(n)}}{\partial \phi_m^j}\frac{\partial \phi_m^j}{\partial x_i}.
\end{eqnarray}


\subsubsection{Details of deformable polygon model}
\label{sec:ch4:DP-details}

In our simulations of the deformable polygon model, each cell is represented by $M=100$ vertices. Enumerating the vertices $\{\mathbf{h}_m\}$ of cell $i$ in clockwise order, the force exerted on vertex $m$ is given by
\begin{equation}
    \mathbf{f}_m=-\frac{\partial E}{\partial \mathbf{h}_m}.
\end{equation}
We need the derivative of area $A_i$ with respect to $\mathbf{h}_m$~\cite{teomy2018confluent}:
\begin{equation}
    \frac{\partial A_i}{\partial\mathbf{h}_m}=\frac{(\mathbf{h}_{m-1}-\mathbf{h}_{m+1})\times\hat{z}}{2},
\end{equation}
and derivatives of $P_m$ and $P_{m-1}$ with respect to $\mathbf{h}_m$~\cite{teomy2018confluent}:
\begin{equation}
    \frac{\partial P_m}{\partial \mathbf{h}_m}=\frac{\mathbf{h}_m-\mathbf{h}_{m+1}}{|\mathbf{h}_m-\mathbf{h}_{m+1}|},\qquad\frac{\partial P_{m-1}}{\partial \mathbf{h}_m}=\frac{\mathbf{h}_m-\mathbf{h}_{m-1}}{|\mathbf{h}_m-\mathbf{h}_{m-1}|},
\end{equation}
where $P_{m}=|\mathbf{h}_m-\mathbf{h}_{m+1}|$, and $\mathbf{h}_{m-1}$ and $\mathbf{h}_{m+1}$ are two adjacent vertices of $\mathbf{h}_m$ in cell $i$.
External forces are applied quasi-statically, starting from zero, and the system is evolved for $5\times10^4$ steps (step size $\Delta t=0.001$) to reach a sequence of steady states.
Vertices are resampled every $1000$ steps by redistributing them along the free-boundary polyline to achieve approximately uniform spacing between adjacent vertices~\cite{weng2022convergent}.
We record the force at which the contact length $P^{(c)}$ falls below the typical segment length between vertices, $\ell_c =2\pi\ell_0/M$, as the detachment force $f_\textrm{detach}^\textrm{DP}$.


\subsection{Analytical details}

To derive the interaction force $f=-\partial E_s/\partial \epsilon$ of a cell doublet [Eq.~\eqref{eq:ch4:force}], we first compute the necessary derivatives. For the angle spanning the non-contacting region $\phi=2\pi-2\operatorname{atan2}\bigl(\sqrt{\ell^2-(\ell-\epsilon)^2},\ell-\epsilon\bigr)$, we obtain
\begin{equation}
    \frac{\partial\phi}{\partial\epsilon}=-\frac{2}{\sqrt{(2\ell-\epsilon)\epsilon}}.
\end{equation}
This gives
\begin{equation}
    \frac{\partial A}{\partial \epsilon} = -2\sqrt{(2\ell-\epsilon)\epsilon},\qquad\frac{\partial P}{\partial \epsilon}=-\frac{2\epsilon}{\sqrt{(2\ell-\epsilon)\epsilon}}.
\end{equation}
We therefore see that the divergent contribution arises from the non-contacting perimeter $P^{(n)}=\ell\phi$ rather than from the area $A$ or the total perimeter $P$.

We can further take the limit $\epsilon\to0^+$ in Eq.~\eqref{eq:ch4:force}, yielding the asymptotic form
\begin{equation}
\label{eq:ch4:force-asymp}
    2\sqrt{2\ell\epsilon}\left[2K_A(\pi\ell^2-A_0) + K_P\frac{(2\pi\ell-P_0)}{\ell}\right] + {\Lambda}\sqrt{\frac{2\ell}{\epsilon}},
\end{equation}
which diverges for $\Lambda\neq 0$, as shown in Fig.~\ref{fig:ch4:doublet}b.
This implies that, with $\Lambda\neq 0$, contacts become increasingly hard to break as their length shrinks to zero in the AFV model.
Note that by using $K_A$ and $L=\ell$ to nondimensionalize terms in the square bracket, we recover the dimensionless criterion for attractive forces given in Ref.~\cite{teomy2018confluent} for the $\Lambda=0$ case:
$
    2(\pi-\tilde{A}_0) + \tilde{K}_P(2\pi-\tilde{P}_0)>0
$.
